\documentclass[a4paper,11pt]{article}
\usepackage[T1]{fontenc}
\usepackage[utf8]{inputenc}
\usepackage{lmodern}
\usepackage{amsmath,amssymb,amsfonts,amsthm}

\newtheorem{claim}{Claim}
\newtheorem{definition}{Definition}

\title{Differential Geometry: Assignment 3}
\author{Aditya Padiyar}
\date{October 2026}

\begin{document}

\maketitle

\begin{definition}
  We call $S \subseteq M$ geodesically convex if for each $p,q \in S$ the following holds:
  \begin{itemize}
    \item There is a geodesic $\gamma:[0,1] \to S$ of $M$ from $p$ to $q$, which is minimizing.
    \item If $\widetilde{\gamma}:[0,1] \to M$ is a minimizing geodesic of $M$ from $p$ to $q$, then $\widetilde{\gamma}=\gamma$
  \end{itemize}
\end{definition}

\begin{definition}
  A loop in $M$ is a piecewise smooth curve $\gamma:[a,b] \to M$ such that $\gamma(a)=\gamma(b)$
\end{definition}

\begin{definition}
  $M$ has trivial holonomy at $p$ if all loops based at $p$ have identity parallel transport
\end{definition}

\begin{definition}
  $M$ has locally trivial holonomy at $p$ if there is a neighborhood $V$ such that $V$ has trivial holonomy at $p$
\end{definition}

\begin{claim}
  If $v,w \in T_pM$ then \[\text{lim}_{t \to 0}\frac{d_M(\exp_p(tv),\exp_p(tw))}{|t|}=||v-w||\]
\end{claim}

\begin{proof}
  Consider the following:
  \begin{itemize}
    \item Let $\gamma_t:[0,1] \to M$ be the minimizing geodesic from $q_1(t)=\exp_p(tv)$ to $q_2(t)=\exp_p(tw)$, for $t \in (-\epsilon,\epsilon)$
    \item We have \[\gamma(s,t)=\exp(q_1(t),s \widetilde{\gamma}(t))\] where $\widetilde{\gamma(t)} \in T_{q_1(t)}M$ is the vector of smallest norm such that \[\exp(q_1(t),\widetilde{\gamma}(t))=q_2(t)\]
    \item The map $(s,t) \to \gamma(s,t)$ is smooth by the inverse function theorem
    \item Let $E(t)=\langle \partial_1 \gamma(0,t) , \partial_1 \gamma(0,t) \rangle $
    \item Then \[E''(t)=2 [ \langle \partial_1 \partial_2     \gamma(0,t),  \partial_1 \partial_2 \gamma(0,t) \rangle + 
                            \langle \partial_1 (\partial_2)^2 \gamma(0,t) , \partial_1 \gamma(0,t)            \rangle ] \]
    \item Now \[(\partial_1^2 \partial_2 \gamma)(s,0)=(D_s)^2 (\partial_2 \gamma)(s,0)=(D_s D_t \partial_1 \gamma)(s,0)\]
          which equals \[(R(\partial_1 \gamma,\partial_2 \gamma) \partial_1 \gamma)(s,0)=0\]
    \item Hence $\partial_2 \gamma (s,0)= sv+(1-s)w$
  \end{itemize}
\end{proof}

\section{Small geodesic balls are geodesically convex}

Consider the following:

\begin{itemize}
  \item Suppose $\exp_p(B_\epsilon) \subseteq W$, where $W$ is a uniformly $\delta$ normal neighborhood of $p$ and $\epsilon < \delta$
  \item Define $g:M \to \mathbb{R}$ by $g(q)=[d_M(q,p)]^2$
  \item Define $h:UTB(\frac{\epsilon}{2},p) \times (-\frac{\epsilon}{2},\frac{\epsilon}{2}) \to M$ by $(q,v,t) \to \exp_q(tv)$
  \item Let $K=[\frac{-\epsilon}{4},\frac{\epsilon}{4}]$
  \item Define $f=g \circ h$
  \item Let $U_0=\{(q,v) \in UTB(\frac{\epsilon}{2},p) :f''(q,v,t)>0 \text{ for all }t \in K\}$
  \item Let $U=\{q \in B(\frac{\epsilon}{2},p): (q,v) \in U_0 \text{ for all } v \in T_qM \text{ with } ||v||=1 \}$
\end{itemize}

\begin{claim}
  $\text{range }h \subseteq W$
\end{claim}

\begin{proof}
  Consider the following:
  \begin{itemize}
    \item Let $q_0=h(q,v,t)$
    \item $d_M(q_0,q)=|t|<\frac{\epsilon}{2}$
    \item $d_M(q_0,p) \leq d_M(q_0,q) + d_M(p,q) < \epsilon$
    \item But we have that $B(\epsilon,p) \subseteq W$
  \end{itemize}
\end{proof}

\begin{claim}
  $g$ is smooth on $B(\epsilon,p)$
\end{claim}

\begin{proof}
  This follows since $g(\exp_p(v))=||v||^2$ and $\exp_p$ is a diffeomorphism from $B_\epsilon$ to $B(\epsilon,p)$
\end{proof}

\begin{claim}
  $f(p,v,t)=t^2$
\end{claim}

\begin{claim}
  We have $f''(a)=\text{Hess}_{h(a)}(g)(h'(a),h'(a))$
\end{claim}

\begin{claim}
  $\text{Hess}_q(g)$ is positive definite for $q \in U$
\end{claim}

\begin{claim}
  $U$ is open.
\end{claim}

\begin{proof}
  Consider the following:
  \begin{itemize}
    \item $U_0$ is open by tube lemma
    \item Choose an orthonormal frame on some neighborhood $V_0$ of $p$. We have the diffeomorphism $\Psi:V_0 \times S^{n-1} \to UTV_0$. 
    \item By tube lemma $U=\{q:(q,v) \in \Psi^{-1}(U_0) \text{ for all }v\}$ is open
  \end{itemize}
\end{proof}

\begin{claim}
    Suppose $q_1,q_2 \in U$ and $\gamma:[0,1] \to U$ is a minimizing geodesic segment from $q_1$ to $q_2$. 
    Then $g \circ \gamma$ is convex
\end{claim}

\begin{proof}
  We have $(g \circ \gamma)''(t)=\text{Hess}_{\gamma(t)}(g)(\gamma'(t),\gamma'(t))>0$
\end{proof}

\begin{claim}
  $V=B(p,\frac{r}{3})$ is geodesically convex, where $B(p,r) \subseteq U$ for some $\delta$
\end{claim}

\begin{proof}
  Suppose $q_1,q_2 \in V$ and $\gamma:[0,1] \to M$ is a minimizing geodesic in $M$ from $q_1$ to $q_2$. Then we have:
  
  \begin{itemize}
    \item $d_M(q_1,\gamma(t))=td_M(q_1,q_2)<\frac{2r}{3}$
    \item $d_M(p,\gamma(t)) \leq d_M(p,q_1) + d_M(q_1,\gamma(t)) < r $
    \item by a previous claim $g \circ \gamma$ is convex
    \item so $d_M(p,\gamma(t)) \leq \text{max}\{d_M(p,q_1),d_M(p,q_2)\}<\frac{r}{3}$
  \end{itemize}

\end{proof}

\section{Riemannian isometries}

Suppose $f:M \to N$ is a surjective metric isometry with respect to the intrinsic metric.
Let $q=f(p)$

\begin{claim}
  There is an open ball $B \subseteq \mathcal{E}_p$ in $T_pM$ centered at the origin and a map $\Psi:B \to \mathcal{E}_q$ such that 
  \begin{itemize}
    \item $f(\exp_p(v))=\exp_q(\Psi(v))$ for all $v \in B$
    \item $\Psi(0)=0$
    \item $\Psi$ is a metric isometry
    \item there is a linear isometry $A:T_pM \to T_qM$ such that $Av=\Psi(v)$ for all $v \in B$
  \end{itemize}
\end{claim}

\begin{proof}
 We have $d(f(\exp_p(v)),q)=d(exp_q(\Psi(v)),q)=d(\exp_p(v),p)$. Hence $||\Psi(v)||=||v||$
\end{proof}

\begin{claim}
  There is a neighborhood $V$ of $p$ such that $f_{|V}$ is smooth
\end{claim}

\begin{claim}
  We have $(Df)_p=A$
\end{claim}

\begin{claim}
  $f$ is a Riemannian isometry
\end{claim}

\section{Flat manifolds}

\begin{claim}
  The following are equivalent
  \begin{enumerate}
    \item If $\gamma_0,\gamma_1:[0,1] \to M$ are smooth curves from $p$ to $q$ that are path homotopic then they have the same parallel transport
    \item The restricted holonomy group $\text{Hol}_0(p)$ is trivial for all $p$
    \item $M$ has locally trivial holonomy at every point
    \item $R=0$ 
  \end{enumerate}
\end{claim}

\begin{proof}
  Consider the following:
  \begin{itemize}
    \item Suppose (1) holds.
          Since the constant loop has identity parallel transport, (2) also holds
    
    \item Suppose (2) holds and let $V$ be a simply connected neighborhood of $p$. 
          Let $\gamma$ be a loop based at $p$ in $V$. 
          Since $\gamma$ is path homotopic to the constant loop, it has identity parallel transport.
          Hence (3) holds.
          
    \item Suppose (3) holds, so that there is a neighborhood $V$ such that $V$ has trivial holonomy at $p$.
          Let $u,v,w \in T_pM$ and define $\Gamma(s,t)=\exp_p(su+tv)$. 
          Applying Theorem 7.11 of Lee, we get $R_p=0$. Hence (4) holds.
  \end{itemize}  
\end{proof}
\end{document}
